\SourceLeaf{306}%
% A544-P0306-D01
in welchem $p$ alle in $2D$ nicht aufgehenden Primzahlen durchlaufen muss, und hieraus folgt
\[
k=\frac{\epsilon\surd D}\pi\Pi\frac1{1-\left(\frac{-D}p\right)\frac1p}
\]

% A544-P0306-D02
Hinsichtlich der Strenge dieser Deduction bleibt aber ein Bedenken übrig, welches sich auf die Methode bezieht, den mittlern Werth der Function $f(m)$ durch successive Elimination aller Primzahlen zu bestimmen; denn wenn es auch einleuchtet, dass der Werth der durch Elimination der ersten $n$ Primzahlen erhaltenen Function $f^{(n)}(m)$ mit dem der Function $f^\infty(m)=1$ übereinstimmt, so lange $m$ kleiner bleibt als die zuletzt eliminirte Primzahl, und dass also durch die Wahl eines hinreichend grossen Werthes $n$ diese Uebereinstimmung bis zu jeder vorher vorgeschriebenen Grösse der Zahl $m$ getrieben werden kann, so ist hiermit allein doch keineswegs erwiesen, dass mit unbegrenzt wachsendem $n$ der mittlere Werth der Function $f^{(n)}(m)$ sich dem mittlern Werthe der Function $f^\infty(m)$, d. h. dem Werthe $1$ unbegrenzt nähert. In welcher Weise der Verf. diese Lücke auszufüllen beabsichtigte, lässt sich aus den vorhandenen Papieren nicht mit Sicherheit erkennen; doch führt die schon oben (in der Note zu I) mitgetheilte Formel
\[
1+4\{A-\frac A3+\frac A5-\frac A7+\frac A9-\frac A{11}\ldots\}
\]
für die Anzahl der Paare von Zahlen, deren Quadratsumme den Werth $A$ nicht übertrifft, zu der Vermuthung, dass der Verf., mit Umgehung des unendlichen Productes, für den mittlern Werth der Function $f(m)$ unmittelbar die unendliche Reihe
\[
\Sigma\left(\frac{-D}n\right)\frac1n
\]
gefunden hat, in welcher $n$ der Grösse nach alle positiven ganzen Zahlen durchlaufen muss, die relative Primzahlen zu $2D$ sind. Die einfachste Art, diesen Uebergang anzudeuten, scheint die folgende zu sein.

% A544-P0306-D03
Ist $\mu$ der grösste aller derjenigen Divisoren einer Zahl $m$, welche relative Primzahlen zu $2D$ sind, und setzt man $\theta(m)=f(\mu)$, so ist $\theta(m)$ diejenige Function, welche durch Elimination aller in $2D$ aufgehenden Primzahlen aus $f(m)$ entsteht, und deren mittlerer Werth nach dem Obigen mit demjenigen der Function $f(m)$ übereinstimmt. Da nun $\theta(m)=\Sigma\left(\frac{-D}n\right)$ ist, wo $n$ alle Divisoren von $\mu$, d. h. alle diejenigen Divisoren von $m$ durchläuft, welche relative Primzahlen zu $2D$ sind, so ergiebt sich die der obigen analoge Formel
\[
\Theta(m)=\theta(1)+\theta(2)+\ldots+\theta(m)=\Sigma\left(\frac{-D}n\right)\frac mn
\]
wo in der Summe rechter Hand der Buchstabe $n$ alle relativen Primzahlen zu $2D$ durchläuft, und von dem Quotienten $\frac mn$ immer nur die grösste in ihm enthaltene ganze Zahl beizubehalten ist. Ordnet man die Glieder dieser Reihe so, dass die Zahlen $n$ ihrer Grösse nach wachsend auf einander folgen, so nimmt der Factor $\frac mn$ fortwährend ab oder doch wenigstens nie zu, und die Reihe bricht ab, sobald $n>m$ wird. Ausserdem ergiebt sich aus dem Fundamentaltheorem in der Theorie der quadratischen Reste und aus der Verallgemeinerung desselben, dass die Summe von je $\varphi(4D)$ auf einander folgenden Werthen des Factors $\left(\frac{-D}n\right)$ verschwindet, woraus folgt, dass die Summe von noch so vielen auf einander folgenden Werthen
\SourceLeaf{307}%
% A544-P0307-D01
desselben ihrem absoluten Werth nach die endliche, nur von dem Determinant $D$ abhängige Grösse $\Delta=\varphi(2D)$ niemals übertrifft. Verbindet man diese beiden Bemerkungen mit einander, so findet man leicht, dass die Summe aller auf das Glied $\left(\frac{-D}n\right)\frac mn$ folgenden Glieder absolut genommen kleiner als $\Delta\frac mn$ ist, und dasss folglich der Quotient $\Theta(m):m$ bei unendlich wachsendem $m$ die in der angegebenen Art geordnete, convergirende unendliche Reihe
\[
\Sigma\left(\frac{-D}n\right)\frac1n
\]
zum Grenzwerthe hat. Nachdem so der gemeinschaftliche mittlere Werth der Functionen $\theta(m)$ und $f(m)$ gefunden ist, erhält man unmittelbar
\[
k=\frac{\epsilon\surd D}\pi\Sigma\left(\frac{-D}n\right)\frac1n
\]

% A544-P0307-D02
Es verdient noch bemerkt zu werden, dass die Artikel 6 und 8 der Abhandlung II auf eine in mancher Beziehung einfachere und auch leicht auszuführende Behandlungsweise des Problems hindeuten, bei welcher nur die Darstellungen ungerader oder sogar nur solcher Zahlen betrachtet werden, die relative Primzahlen zu $2D$ sind.
\BookRule{35mm}

% A544-P0307-D03
\pdfbookmark[2]{Zu VI und VII.}{nexus-remarks-vi-vii}
\begin{center}\large Zu VI und VII.\end{center}

% A544-P0307-D04
Die Art, wie der Verf. die Summation der Reihe $\Sigma\left(\frac{-D}n\right)\frac1n$ ausgeführt hat, ergiebt sich aus einigen speciellen Beispielen, welche sich auf einzelnen Blättern vorfinden.

% A544-P0307-D05
Ist $D\equiv3\pmod4$, so folgt aus dem Fundamentaltheorem in der Theorie der quadratischen Reste mit Benutzung der Reihe
\[
\operatorname{cotang}u=\frac1u+\frac1{u-\pi}+\frac1{u+\pi}+\frac1{u-2\pi}+\frac1{u+2\pi}+\ldots
\]
dass
\[
\Sigma\left(\frac{-D}n\right)\frac1n=\Sigma\left(\frac nD\right)\frac1n=\frac\pi{2D}\Sigma\left(\frac\nu D\right)\operatorname{cotang}\frac{\nu\pi}{2D}
\]
ist, wo $\nu$ alle relativen Primzahlen zu $2D$ durchläuft, die kleiner als $D$ sind; setzt man
\[
\surd-1=i,\quad\cos\frac{2\pi}D+i\sin\frac{2\pi}D=r
\]
und bezeichnet mit $\mu$ alle relativen Primzahlen zu $D$, welche nicht grösser als $D$ sind, so lässt sich die vorstehende Summe leicht in die folgende umformen
\[
\Sigma\left(\frac{-D}n\right)\frac1n=\frac{\pi i}{4D}\left(\frac2D\right)\Sigma\left(\frac\mu D\right)\frac{r^\mu-1}{r^\mu+1}
\]
Wendet man nun die für jede Wurzel $\omega$ der Gleichung $\omega^D=1$ gültige Formel
\SourceLeaf{308}%
% A544-P0308-D01
\[
\frac{\omega-1}{\omega+1}=\Sigma(-1)^{\alpha-1}\omega^\alpha
\]
an, in welcher $\alpha$ die Zahlen $1,2,3\ldots(D-1)$ durchlaufen muss, so erhält man durch Umkehrung der Summationsordnung
\[
\Sigma\left(\frac{-D}n\right)\frac1n=\frac{\pi i}{4D}\left(\frac2D\right)\Sigma(-1)^{\alpha-1}\Sigma\left(\frac\mu D\right)r^{\alpha\mu}
\]
Die auf $\mu$ bezügliche Summation lässt sich bekanntlich mit Hülfe der in der Abhandlung \emph{Summatio quarundam serierum singularium} bewiesenen Sätze ausführen; beschränkt man sich auf den Fall, in welchem $D$ durch kein Quadrat theilbar ist, so findet man allgemein
\[
\Sigma\left(\frac\mu D\right)r^{\alpha\mu}=\left(\frac\alpha D\right)i^{\left(\frac{D-1}2\right)^2}\surd D
\]
wo $\left(\frac\alpha D\right)=0$ gesetzt werden muss, falls $\alpha$ keine relative Primzahl zu $D$ ist. In dem Fall $D\equiv3\pmod4$ erhält man daher
\[
\Sigma\left(\frac{-D}n\right)\frac1n=\frac\pi{4\surd D}\left(\frac2D\right)\Sigma(-1)^\alpha\left(\frac\alpha D\right)=\frac\pi{2\surd D}\Sigma\left(\frac{\alpha^{\prime}}D\right)
\]
wo $\alpha^{\prime}$ alle relativen Primzahlen zu $D$ durchläuft, die kleiner als $\frac12D$ sind; da endlich $\epsilon=2$ ist, so wird die Anzahl der Classen
\[
k=\Sigma\left(\frac{\alpha^{\prime}}D\right)
\]

% A544-P0308-D02
Ist dagegen $D\equiv1\pmod4$, so erhält man mit Benutzung der Reihe
\[
\operatorname{cosec}u=\frac1u-\frac1{u-\pi}-\frac1{u+\pi}+\frac1{u-2\pi}+\frac1{u+2\pi}-\ldots
\]
auf ähnliche Weise
\[
\begin{aligned}
\Sigma\left(\frac{-D}n\right)\frac1n&=\Sigma(-1)^{\frac{n-1}2}\left(\frac nD\right)\frac1n
=\frac\pi{2D}\Sigma(-1)^{\frac{\nu-1}2}\left(\frac\nu D\right)\operatorname{cosec}\frac{\nu\pi}{2D}\\
&=\frac\pi{2D}\Sigma\left(\frac\mu D\right)\frac{r^\mu}{r^{2\mu}+1}
\end{aligned}
\]
wo die Buchstaben $\nu$ und $\mu$ die frühere Bedeutung haben; schliesst man den evidenten Fall $D=1$ aus und wendet die für jede Wurzel $\omega$ der Gleichung $\omega^D=1$ (mit Ausnahme von $\omega=1$) gültige Formel
\[
\frac\omega{\omega\omega+1}=1+\Sigma\omega^{4\alpha^{\prime\prime}}+\Sigma\omega^{D-4\alpha^{\prime\prime}}
\]
an, in welcher $\alpha^{\prime\prime}$ die Zahlen $1,2,3\ldots\frac14(D-1)$ durchlaufen muss, so ergiebt sich, wieder unter der Beschränkung, dass $D$ durch kein Quadrat theilbar ist,
\[
\Sigma\left(\frac{-D}n\right)\frac1n=\frac\pi{\surd D}\Sigma\left(\frac{\alpha^{\prime\prime}}D\right)
\]
und hieraus, da $\epsilon=2$ ist,
\[
k=2\Sigma\left(\frac{\alpha^{\prime\prime}}D\right)
\]

\SourceLeaf{309}%
% A544-P0309-D01
Ganz ähnlich würden sich die Fälle behandeln lassen, in welchen $D$ gerade ist. —

% A544-P0309-D02
Was die Bestimmung der Classen-Anzahl für \emph{positive} Determinanten $D$ betrifft, so finden sich ausser der im Text mitgetheilten Schlussformel nur einzelne geometrische Figuren vor, welche Hyperbel-Sectoren von endlichen Dimensionen darstellen, und neben denselben Ungleichungen, durch welche die Punkte, deren Coordinaten die Variabeln der quadratischen Formen sind, in das Innere eines solchen Hyperbel-Sectors gedrängt werden. Diese Hyperbel-Sectoren treten an die Stelle der Ellipsen, welche den quadratischen Formen von negativen Determinanten entsprechen, und durch die Bestimmung ihres Flächeninhalts ergiebt sich wieder die mittlere Darstellungsanzahl, wenn nämlich nur solche Darstellungen zugelassen werden, bei welchen die Variabeln den eben erwähnten Ungleichungen Genüge leisten. Andererseits dienen diese Ungleichungen dazu, aus den unendlich vielen Darstellungen einer Zahl $m$, welche alle zu einer und derselben Wurzel $n$ der Congruenz $nn-D\equiv0\pmod{\frac m{\mu\mu}}$ gehören und welche den sämmtlichen Auflösungen der Gleichung $tt-Duu=1$ entsprechen (vergl. Disqq. Arithm. art. 205), eine einzige zu isoliren und alle andern auszuschliessen. Die Anzahl aller zugelassenen Darstellungen der Zahl $m$ durch den Complex aller nicht eigentlich äquivalenten formae proprie primitivae ist dann gleich dem Werth der Function $f(m)$, in welcher nur $-D$ durch $D$ zu ersetzen ist, und aus der Betrachtung der Eigenschaften derselben ergiebt sich, wie früher bei negativen Determinanten, ein zweiter Ausdruck für die mittlere Darstellungsanzahl; die Vergleichung desselben mit dem vorher durch geometrische Betrachtungen abgeleiteten Werthe führt dann unmittelbar zu der Bestimmung der Anzahl der Classen.
\BookRule{35mm}

% A544-P0309-D03
\pdfbookmark[2]{Zu VIII.}{nexus-remarks-viii}
\begin{center}\large Zu VIII.\end{center}

% A544-P0309-D04
Hier bedeutet $p$ eine positive Primzahl von der Form $4n+1$; die Bezeichnung stimmt mit der in der Abhandlung \emph{Theoria residuorum biquadraticorum} I. art. 23 angewendeten überein; es ist also
\[
\begin{gathered}
f\equiv1\pdot2\pdot3\dotsfour\tfrac12(p-3)\pdot\tfrac12(p-1)\pmod p\\
p=aa+bb;\quad a\equiv1\pmod4;\quad b\equiv af\pmod p
\end{gathered}
\]
die mit $\alpha,\beta$ bezeichneten Zahlen sind durch die Zerlegung $p=\alpha\alpha+2\beta\beta$ bestimmt. Die Columne $f$ ist den beiden vorgefundenen Tabellen hinzugefügt; ausserdem sind einige Lücken in denselben ausgefüllt.

% A544-P0309-D05
Der im Text aufgestellte Satz hängt mit dem biquadratischen Charakter der Zahl $2$ zusammen; da nämlich (vergl. \emph{Theoria resid. biqu.} I. art. 21)
\[
2^{\frac{p-1}4}\equiv f^{\frac12 b}\pmod p
\]
ist, so folgt aus der Congruenz
\[
b\equiv2m+a-1\pmod8
\]
die andere
\SourceLeaf{310}%
% A544-P0310-D01
\[
2^{\frac{p-1}4}\equiv f^{m+\frac{a-1}2}\pmod p
\]
und umgekehrt jene aus dieser. Der Beweis dieser letztern Congruenz ergiebt sich leicht auf folgende Art. Ist $\mu$ die Anzahl der quadratischen Reste $\alpha_1$, welche zwischen $0$ und $\frac14p$ liegen, so ist (nach VII)
\[
m=2\mu-\tfrac14(p-1)
\]
und die Anzahl der quadratischen Reste $\alpha_2$, welche zwischen $\frac14p$ und $\frac12p$ liegen, ist $=\frac14(p-1)-\mu$. Ist nun $p\equiv1\pmod8$, also die Zahl $2$ quadratischer Rest, so stimmen die Zahlen $2\alpha_1$ und $p-2\alpha_2$ im Complex mit den Zahlen $\alpha_1$ und $\alpha_2$ überein, und bezeichnet man das Product dieser Zahlen mit $A$, so ergiebt sich
\[
2^{\frac{p-1}4}A\equiv(-1)^{\frac14(p-1)-\mu}A\pmod p
\]
und folglich
\[
2^{\frac{p-1}4}\equiv(-1)^\mu\equiv f^{2\mu}\equiv f^{m+\frac14(p-1)}\pmod p
\]
da ferner in diesem Fall $b\equiv0\pmod4$, und folglich
\[
\frac{p-1}4=(a+1)\frac{a-1}4+\frac{bb}4\equiv2\frac{a-1}4\equiv\frac{a-1}2\pmod4
\]
ist, so erhält man die zu beweisende Congruenz
\[
2^{\frac{p-1}4}\equiv f^{m+\frac{a-1}2}\pmod p
\]

% A544-P0310-D02
Ist dagegen $p\equiv5\pmod8$, also die Zahl $2$ quadratischer Nichtrest, so stimmen die Zahlen $2\alpha_1$ und $p-2\alpha_2$ mit den sämmtlichen zwischen $0$ und $\frac12p$ liegenden quadratischen Nichtresten überein; bezeichnet man ihr Product mit $B$, und das Product der Zahlen $\alpha_1$ und $\alpha_2$ wieder mit $A$, so ist
\[
f\equiv AB,\quad(-1)^{\frac{p-1}4-\mu}2^{\frac{p-1}4}A\equiv B\pmod p
\]
erhebt man diese beiden Congruenzen zum Quadrat, indem man berücksichtigt, dass
\[
ff\equiv-1,\quad2^{\frac{p-1}2}\equiv-1\pmod p
\]
ist, so erhält man
\[
-1\equiv AABB,\quad-AA\equiv BB
\]
und hieraus $A^4\equiv+1$; da nun $A$ ein Product aus quadratischen Resten, also $AA$ ein Product aus biquadratischen Resten und folglich selbst ein biquadratischer Rest ist, so muss $AA\equiv+1$ sein, weil $-1$ ein biquadratischer Nichtrest ist. Hieraus folgt
